\chapter{Adquisición distribuida y comunicación serial síncrona}
\section{\obj}
\capacidad
\begin{itemize}
\item Transmitir datos mediante un protocolo serial síncrono de baja velocidad desde un NI myDAQ maestro.
\item Recibir, visualizar y registrar los datos transmitidos en un NI myDAQ esclavo.
\end{itemize}

\section{\mat}
\textbf{A suministrar por la Escuela:}
\begin{itemize}
\item 2 dispositivos de adquisición de datos NI myDAQ
\item 1 termistor TDC310 de \SI{10}{\kilo\ohm}
\item 1 resistor de \SI{10}{\kilo\ohm}
\item 1 capacitor de \SI{0.1}{\micro\farad}
\end{itemize}

\textbf{A suministrar por el estudiante:}
\begin{itemize}
\item 1 breadboard
\item Cables de interconexión (jumpers) macho-macho
\end{itemize}

\section{\pro}
\begin{enumerate}
\item Implemente un sistema distribuido compuesto por dos estaciones:
    \begin{itemize}
        \item una estación emisora, basada en un NI myDAQ, encargada de medir la señal del sensor de temperatura y transmitir el dato
        \item una estación receptora, basada en un segundo NI myDAQ, encargada de recibir, visualizar y registrar el dato transmitido
    \end{itemize}
\item Conecte ambas estaciones mediante un enlace serial síncrono de baja velocidad tal y como se indica en la Figura \ref{fig:L4F1}. Utilice la siguiente asignación de líneas digitales:
    \begin{itemize}
        \item DIO0: reloj (\emph{CLK})
        \item DIO1: datos (\emph{DATA})
        \item DIO2: sincronismo (\emph{SYNC})
        \item DGND: tierra (\emph{GND})
    \end{itemize}
\begin{figure}[H]
    \centering
    \begin{circuitikz}
        \draw[dashed,blue]
        (-4,5) rectangle (-0.5,-3.5)
        node[midway]{}
        ;
        \draw[dashed,blue]
        (1.5,5) rectangle (5,-3.5)
        node[midway]{}
        ;

        \draw
        (-2.25,4.2) node{\small myDAQ emisor}
        (3.25,4.2) node{\small myDAQ receptor}
        (-3.4,3.2) node[left]{DIO0}
        (-3.4,1.6) node[left]{DIO1}
        (-3.4,0.0) node[left]{DIO2}
        (-3.4,-1.6) node[left]{DGND}
        (2.1,3.2) node[left]{DIO0}
        (2.1,1.6) node[left]{DIO1}
        (2.1,0.0) node[left]{DIO2}
        (2.1,-1.6) node[left]{DGND}
        ;

        \draw[red, thick, -latex]
        (-0.5,3.2) -- (1.5,3.2)
        node[midway, above]{CLK}
        ;
        \draw[orange, thick, -latex]
        (-0.5,1.6) -- (1.5,1.6)
        node[midway, above]{DATA}
        ;
        \draw[green!60!black, thick, -latex]
        (-0.5,0.0) -- (1.5,0.0)
        node[midway, above]{SYNC}
        ;
        \draw[black, thick]
        (-0.5,-1.6) -- (1.5,-1.6)
        node[midway, below]{GND}
        ;
    \end{circuitikz}
    \caption{Enlace serial síncrono entre dos NI myDAQ usando DIO0 como reloj, DIO1 como datos y DIO2 como sincronismo.}
    \label{fig:L4F1}
\end{figure}
\item En la estación emisora, cree un instrumento virtual (\emph{.vi}) que realice lo siguiente:
    \begin{itemize}
        \item Leer la señal analógica del sensor de temperatura cada \SI{1000}{\milli\second}
        \item Convertir la señal medida a una temperatura en unidades de \si{\degreeCelsius}
        \item Cuantizar el valor de temperatura a un número entero adecuado para su transmisión
        \item Formar una trama digital y transmitirla bit a bit utilizando las líneas \emph{SYNC}, \emph{CLK} y \emph{DATA}
    \end{itemize}
\item En la estación receptora, cree un instrumento virtual (\emph{.vi}) que realice lo siguiente:
    \begin{itemize}
        \item Detectar el inicio de la trama mediante la línea \emph{SYNC}
        \item Leer cada bit recibido usando la línea de reloj
        \item Reconstruir el valor numérico transmitido
        \item Convertir el valor recibido a temperatura y mostrarlo en el panel frontal
        \item Graficar la temperatura recibida en tiempo real
        \item Guardar los datos recibidos en un archivo separado por comas (\emph{.csv}), de forma que la primera columna sea el tiempo y la segunda la temperatura recibida
    \end{itemize}
\item Utilice una frecuencia de transmisión suficientemente baja para garantizar la correcta operación del enlace serial síncrono implementado por software.
\item Verifique experimentalmente que el valor transmitido por la estación emisora coincide con el valor reconstruido por la estación receptora.
\item Incluya en su informe lo siguiente:
    \begin{itemize}
        \item Un diagrama del enlace entre ambos NI myDAQ indicando las líneas \emph{SYNC}, \emph{CLK}, \emph{DATA} y \emph{GND}
        \item Una captura de pantalla del diagrama de bloques del emisor
        \item Una captura de pantalla del diagrama de bloques del receptor
        \item Una captura de pantalla del panel frontal del receptor
        \item Una gráfica de la temperatura recibida a partir del archivo \emph{.csv}, usando Python. Incluya títulos de ejes adecuados
        \item Una breve comparación entre los valores transmitidos y los valores recibidos
    \end{itemize}
\end{enumerate}
